\chapter{Introducción}
\label{chap:introduccion}

Este Trabajo de Fin de Grado aborda el análisis automático de electrocardiogramas representados como imágenes. La entrada corresponde al trazado ya renderizado. Por tanto, el análisis se concentra en la morfología que conserva esa representación.

El caso de estudio es la sospecha de síndrome de Brugada, una canalopatía hereditaria poco frecuente asociada al riesgo de arritmias ventriculares y muerte súbita. La escasez de pacientes confirmados dificulta el entrenamiento y la evaluación de modelos con los volúmenes de datos habituales en aprendizaje profundo.

La pregunta metodológica es cómo analizar ECG como imágenes cuando solo se dispone de unos pocos pacientes etiquetados. La representación visual conserva información morfológica, pero prescinde de parte de la estructura de la señal original.

Para responder a esta pregunta se comparan dos formas de aprovechar el mismo conjunto reducido de pacientes etiquetados. En la primera, una red neuronal convolucional (CNN) actualiza sus pesos a partir de \(K\) pacientes de entrenamiento. En la segunda, un modelo de visión y lenguaje (VLM) mantiene sus pesos fijos y recibe \(K\) demostraciones dentro del contexto. Por tanto, \(K\) expresa el número de ejemplos disponibles para cada método, aunque su función sea distinta en el entrenamiento supervisado y en el aprendizaje en contexto.

\section{Motivación clínica}

El síndrome de Brugada es un trastorno arrítmico hereditario relacionado con alteraciones de los canales de sodio cardíacos. Predispone a fibrilación ventricular y muerte súbita, especialmente en varones jóvenes durante el reposo o el sueño \cite{zeppenfeld2022,wilde2023scbt}. Su prevalencia estimada se sitúa entre una de cada 2\,000 y una de cada 5\,000 personas. Sin embargo, es probable que esté infradiagnosticado porque muchos portadores permanecen asintomáticos.

El patrón electrocardiográfico tipo~1 se caracteriza por una elevación del segmento ST de morfología convexa de al menos 2\,mm en una derivación precordial derecha, seguida de una onda T negativa \cite{wilde2023scbt}. Los patrones tipo~2 y tipo~3 no son diagnósticos por sí solos. También existen fenocopias producidas por fiebre, alteraciones electrolíticas, isquemia o fármacos. Estas pueden imitar el trazado y deben excluirse antes de confirmar la enfermedad \cite{sieira2017}.

La dificultad aumenta porque el patrón tipo~1 puede ser intermitente. Un ECG basal normal no descarta el síndrome y un trazado sospechoso debe interpretarse dentro del contexto clínico. Esa variabilidad convierte el síndrome de Brugada en un buen caso para estudiar métodos de apoyo a la revisión del ECG, ya que obliga a distinguir entre reconocer un patrón visual y establecer una conclusión clínica.

\subsection{Papel del sistema propuesto}

El sistema se evalúa como posible apoyo previo a la confirmación diagnóstica, limitado a priorizar la revisión de trazados sospechosos. En este marco, una predicción positiva indica que el ECG debe revisarse con mayor atención dentro del circuito clínico habitual. La valoración especializada, la indicación de pruebas complementarias y las decisiones terapéuticas quedan fuera del alcance del modelo.

Este alcance condiciona la interpretación de las métricas. En una enfermedad de baja prevalencia, un clasificador con sensibilidad alta puede generar muchas alertas no confirmadas si su especificidad es insuficiente. Aumentar la sensibilidad reduce el riesgo de retrasar la revisión de un caso sospechoso, pero suele incrementar los falsos positivos y la carga sobre el circuito asistencial. La evaluación debe considerar, por tanto, tanto la capacidad de detectar positivos como la proporción de pacientes normales derivados de forma innecesaria.

\section{Motivación metodológica}

Los modelos profundos han alcanzado un buen rendimiento en el análisis de ECG cuando disponen de grandes colecciones de pacientes y etiquetas consistentes \cite{ribeiro2020,sangha2022}. Estas condiciones rara vez se cumplen en enfermedades de baja prevalencia. En el síndrome de Brugada, cada caso confirmado exige revisión especializada y, en muchos pacientes, información clínica adicional o una prueba farmacológica.

La escasez de etiquetas coincide además con una diferencia de dominio. Las imágenes sintéticas permiten controlar la morfología. La heterogeneidad de los ECG clínicos y la equivalencia entre las referencias quedan fuera de ese control. El problema experimental consiste en determinar qué parte del aprendizaje con pocos casos se conserva al cambiar de dominio y de referencia.

Sobre esta base, el trabajo aborda tres preguntas:

\begin{enumerate}
  \item ¿Cómo cambia el rendimiento al aumentar el número de pacientes etiquetados en un dominio sintético controlado?
  \item ¿En qué medida se conserva ese rendimiento al evaluar ECG clínicos del HUCA?
  \item ¿Cómo influye en el rendimiento del VLM que las demostraciones procedan del simulador o de la muestra clínica?
\end{enumerate}

\section{Alcance y contribuciones}

El estudio se limita a un análisis experimental y retrospectivo de métodos de aprendizaje con pocos datos etiquetados. Sus resultados no validan un producto sanitario ni permiten estimar su utilidad clínica prospectiva. La comparación utiliza dos tipos de referencia. En el dominio sintético, las etiquetas se derivan de reglas aplicadas a tres hallazgos morfológicos del trazado. En la muestra del Hospital Universitario Central de Asturias (HUCA), la etiqueta se asigna por paciente y distingue entre casos confirmados de síndrome de Brugada y pacientes normales.

El trabajo aporta un protocolo para estudiar ECG como imágenes, un simulador QRS/ST revisado y el procesamiento de 317 pacientes del HUCA. También establece una referencia basada en CNN para los dominios sintético y real y evalúa Gemma~4 en zero-shot, ICL estándar e ICL balanceado.

\section{Estructura de la memoria}

La memoria se organiza desde el contexto del problema hasta la interpretación de los resultados. El capítulo~\ref{chap:relacionados} sitúa el trabajo frente al estado del arte. El capítulo~\ref{chap:fundamentos} introduce los fundamentos clínicos y técnicos. El capítulo~\ref{chap:diseno} formula el protocolo común. El capítulo~\ref{chap:simulador} describe el simulador. El capítulo~\ref{chap:modelos} detalla la configuración de los modelos.

Dentro del capítulo de modelado, las secciones~\ref{chap:cnn_enfoque} y \ref{chap:vlm} presentan, respectivamente, la configuración de la CNN y la del enfoque VLM/ICL. El capítulo~\ref{chap:resultados} presenta los hallazgos en el simulador y en los ECG reales. El capítulo~\ref{chap:conclusiones} los discute frente a la literatura, delimita sus limitaciones y formula las conclusiones. El capítulo~\ref{chap:gestion_etica} recoge la gestión ética, regulatoria y de privacidad.
